% Chapter 2 -- Rubric: PO1 (1.1), PO2 (research literature), feeds PO3 and PO11
\chapter{Preliminary Research}
\label{ch:research}

\section{Literature Review}
\label{sec:literature}

\subsection{Modelling what a learner knows}

The problem of inferring an unobservable knowledge state from observable
responses has a long history in intelligent tutoring systems. Bayesian
Knowledge Tracing (BKT), introduced by Corbett and Anderson
\cite{corbett1995knowledge}, remains the reference formulation: a skill's
mastery is a hidden binary variable, and four parameters govern its evolution
--- a prior $P(L_0)$, a learning rate $P(T)$, and two noise terms, the slip
probability $P(S)$ (knowing but answering wrongly) and the guess probability
$P(G)$ (not knowing but answering correctly). The model's durability comes from
the fact that these four parameters are interpretable to a teacher, which is not
true of most of its successors.

The literature has extended BKT along three axes that are directly relevant
here. The first is \emph{parameter individualisation}: Yudelson et al.
\cite{yudelson2013individualized} fit per-student parameters and report
improved predictive accuracy, establishing that treating all learners as
identical is a real source of error. The second is \emph{contextualisation of
noise}: Baker et al. \cite{baker2008more} estimate slip and guess per response
rather than per skill, on the argument that a guess on a four-option item is not
the same event as a guess on an open response. Pardos and Heffernan
\cite{pardos2011kt} make the related argument for item difficulty. The third is
\emph{structure between skills}: Käser et al. \cite{kaser2017dynamic} model
students with dynamic Bayesian networks in which skills are linked by
prerequisite relations, and report that exploiting that structure improves
prediction over independent per-skill models.

Competing model families exist. Learning Factors Analysis \cite{cen2006learning}
and Performance Factors Analysis \cite{pavlik2009performance} replace the
hidden-state formulation with logistic regression over counts of prior successes
and failures. Deep Knowledge Tracing \cite{piech2015deep} replaces it with a
recurrent neural network and reports substantial accuracy gains on large
datasets. These are relevant as alternatives considered
(\Cref{sec:alternatives}) but share a property that disqualified them here:
their internal state is not interpretable as ``this student has not mastered
this specific skill,'' which is exactly the output this system must produce in
order to act on it.

A separate and older tradition, knowledge space theory
\cite{doignon1985spaces,falmagne2011learning}, formalises a domain as a
partially ordered set of knowledge states and asks which questions most
efficiently locate a learner within it. This is the theoretical ancestor of the
prerequisite graph used here, and it supplies the central insight that a
well-structured domain permits assessment far shorter than the number of skills,
because observing one response constrains many states at once.

\subsection{Cognitive level as a modelling variable}

Bloom's taxonomy \cite{bloom1956taxonomy}, in its revised form
\cite{anderson2001revised}, classifies learning objectives into six ascending
cognitive levels from Remember to Create. Its use in tutoring systems is usually
editorial --- a way of labelling items --- rather than a variable inside the
student model. The observation this project builds on is that the taxonomy has a
direct probabilistic reading: a four-option recall question can be answered
correctly by chance far more easily than an analysis question, and a learner who
genuinely knows an analysis-level skill is more likely to slip on it through a
procedural error. That makes the Bloom level a natural conditioning variable for
$P(G)$ and $P(S)$, which is the refinement Baker et al. \cite{baker2008more}
argue for on empirical grounds but implemented from a pedagogical prior rather
than a fitted one.

Vygotsky's zone of proximal development \cite{vygotsky1978mind} supplies the
selection rule: instruction is most effective just beyond what the learner can
already do unaided. Translated into this system's terms, the next question
should sit one Bloom level above the learner's current estimated band --- which
is the rule the diagnostic implements.

\subsection{Automatic generation of assessment items}

Kurdi et al. \cite{kurdi2020systematic} survey automatic question generation for
education and report the field's persistent difficulty: generated items are
often syntactically fine but pedagogically hollow, and evaluation of their
quality is usually thin. Retrieval-augmented generation
\cite{lewis2020retrieval} addresses part of this by grounding generation in
retrieved source material rather than in model memory, which both improves
factual accuracy and ties items to a specific curriculum. Multilingual sentence
embeddings \cite{reimers2019sentence} make such retrieval feasible over a
non-English corpus.

\subsection{What the literature establishes, and what it leaves open}

Established: BKT is an adequate and interpretable model of per-skill mastery;
its noise parameters should not be uniform across items; prerequisite structure
between skills carries real predictive value; and short adaptive assessment is
possible when the domain is well structured.

Left open, and therefore the space this project works in: the literature
provides no standard method for \emph{constructing} a fine-grained prerequisite
graph for a specific national syllabus --- the graphs in published work are
either small, hand-built by domain experts, or learned from response data that
does not exist before a system is deployed. Nor does it settle how a
Bloom-conditioned, graph-propagating variant of BKT should behave when the
question bank is incomplete, which is the normal condition of a system built
before it has users.

\section{Existing Solutions and Technologies}
\label{sec:related}

\begin{table}[htbp]
  \centering
  \caption{Existing solutions surveyed, and their limitations for this problem}
  \label{tab:existing}
  \begin{tabularx}{\textwidth}{>{\hsize=0.6\hsize}Y>{\hsize=0.75\hsize}Y>{\hsize=1.65\hsize}Y}
    \toprule
    Category & Examples & Strengths and limitations for this problem \\
    \midrule
    Adaptive commercial platforms & ALEKS; Knewton-style adaptive engines &
    Strength: mature adaptive engines with knowledge-space foundations and
    demonstrated scale. Limitation: closed, licensed per seat, and built around
    Western curricula --- no HSC syllabus coverage, no Bangla content, and no
    ability to add either. \\
    Bangladeshi ed-tech platforms & 10 Minute School, Shikho and similar &
    Strength: strong local content, Bangla-first, priced for the local market,
    already trusted by the target cohort. Limitation: predominantly
    content-delivery and fixed mock tests; assessment reports a score per test
    or per chapter rather than inferring a per-skill knowledge state, so the
    student still has to decide what to study next. \\
    Generic quiz and spaced-repetition tools & Quizlet, Anki, Google Forms &
    Strength: free, simple, and effective for pure recall
    \cite{ebbinghaus1885memory}. Limitation: no model of prerequisite structure
    at all --- a card is either due or not due --- so they cannot diagnose
    \emph{why} an item was failed or redirect to an upstream deficit. \\
    Printed question banks and coaching centres & The incumbent method &
    Strength: authoritative content, human explanation, peer environment.
    Limitation: fixed pace for a heterogeneous room; diagnosis at the
    granularity of a chapter, not a skill; and the marginal cost of one-to-one
    attention is precisely what prevents it scaling
    \cite{vanlehn2011relative}. \\
    Open-source ITS frameworks & pyBKT and similar research toolkits &
    Strength: reference implementations of the model family, useful for
    validation. Limitation: libraries, not products --- they supply the
    inference but neither the ontology, the content, nor the serving layer,
    which is where the majority of this project's effort actually went. \\
    \bottomrule
  \end{tabularx}
\end{table}

\noindent\textbf{The gap.} No surveyed system combines (i) a fine-grained
prerequisite model of \emph{this} syllabus, (ii) inference that is interpretable
per skill and propagates across prerequisite structure, and (iii) Bangla content
with correct mathematical typesetting. Platforms that have the content lack the
model; platforms that have the model lack the content and the language. ATLAS
targets that intersection.

\section{Market and User Research}
\label{sec:market}

\textbf{User base.} The primary users are HSC-completing students in Bangladesh
preparing for engineering admission tests, typically in a concentrated period
between the HSC examination and the admission tests. Secondary users are
repeat candidates and students in the HSC years using the system for syllabus
practice rather than test preparation.

\textbf{What users do instead.} Current practice, established from the team's
own experience of the same examination and from the structure of the material
that dominates the market, is: solve past-paper compilations; attend coaching
classes; sit periodic mock tests that return a rank and a score. The gap this
leaves is diagnostic, not informational --- students do not lack content, they
lack a reliable answer to ``what should I study next?''

\textbf{Demand characteristics.} Willingness to spend on preparation is
demonstrably high, since coaching enrolment is close to universal among serious
candidates, but willingness to spend on \emph{software} is low; the realistic
price point for a digital tool is at or near zero for the student, with revenue
arriving through institutions if at all (\Cref{sec:economics}).

\textbf{Evidence and its limits.} The characterisation above is grounded in the
team's direct familiarity with the examination and an informal survey of
available platforms and printed material. It is \emph{not} grounded in a
structured survey of prospective users with a defined sample; no such study was
conducted. This is a genuine weakness of the project's user research, stated
here rather than obscured, and \Cref{sec:conclusion} lists it as future work.
% TODO (optional but recommended): if the team can run even a small structured
% survey of 30--50 candidates before submission, replace the paragraph above
% with real figures -- it converts the weakest section of this chapter into a
% strength.

\section{Feasibility Study}
\label{sec:feasibility}

\textbf{Technical.} Feasible, and demonstrated. The inference model is
computationally trivial (a constant-time update per response plus an ancestor
traversal bounded by graph depth, measured at 9). The content pipeline was the
real technical risk: it depends on a large language model producing correct
Bangla mathematical content at scale. That risk was retired empirically --- an
initial attempt using local models on CPU hardware proved infeasible on
throughput grounds and was abandoned, and a hosted API with a
generate-then-verify design was substituted successfully
(\Cref{sec:generation-pipeline}).

\textbf{Economic.} Feasible at prototype scale at essentially zero marginal
cost: every service used (database, model API, hosting) has a free tier adequate
for development and demonstration. It is \emph{not} feasible to assume those
free tiers extend to production; \Cref{sec:economics} models the real operating
cost.

\textbf{Operational.} The adoption question is the binding one. The system asks
a student to spend 30 questions on a diagnostic before it gives them anything,
which is a real cost to a user under time pressure who is accustomed to
immediate practice. Mitigation is to make the diagnostic itself feel like
practice --- it serves genuine past-paper-style questions, not artificial
calibration items.

\textbf{Legal and regulatory.} Feasible with attention. The system stores
learner performance data, which is personal data; the source content is drawn
from copyrighted textbooks and question banks, which constrains redistribution.
Both are addressed in \Cref{sec:standards} and \Cref{sec:ip}.

\textbf{Conclusion.} The project is feasible, and the principal assumption it
depends on is that a machine-constructed skill ontology, validated structurally
but not yet by domain experts, is an adequate foundation for the inference built
on top of it. That assumption is explicitly flagged as unretired
(\Cref{sec:eval-limits}).
